% !TEX root = ./jfla2027_arbres_normatifs.tex

\section{Introduction}\label{sec:introduction}

\subsection{Le raisonnement du juriste} \label{sec:juriste}
Supposons qu'une salariée demande quelle indemnité lui serait due en cas de
licenciement, ou souhaite connaître la durée de son préavis. Même si toutes les
formules candidates qui représentent fidèlement le droit français sont
exécutables, la réponse ne peut se réduire à un nombre, ni même à une fonction.
La question préalable est : \emph{quelle règle juridique est applicable à cette
salariée, pour cette question, à la date juridiquement pertinente ?}

Un juriste saisi de cette demande ne calcule qu'en dernier lieu. Il commence par
\emph{qualifier} les faits et la demande : la rupture est-elle un licenciement,
pour quel motif, et quel avantage la salariée réclame-t-elle exactement ? Nous
appelons \emph{matière} la question juridique ainsi précisée. Il recense
ensuite les \emph{sources} qui peuvent régir cette matière : la loi, ici le
Code du travail ; la convention collective de branche dont relève l'entreprise,
accord négocié entre syndicats de salariés et organisations d'employeurs d'un
secteur d'activité ; un éventuel accord d'entreprise ; le contrat de travail.
Pour chaque source, il retient la \emph{version} du texte en vigueur à la date
pertinente, ici celle de la notification du licenciement, et vérifie que la
salariée en remplit les \emph{conditions d'application}. Si plusieurs textes
subsistent, il règle leur \emph{concours}, au besoin en calculant la
réponse de chacun. 

\subsection{Vocabulaire juridique}\label{sec:vocabulaire}
%
Plusieurs notions de droit, employées dans cet article, doivent être introduites.

\paragraph{Hiérarchie des normes.} Les sources du droit sont ordonnées par
  rang : la Constitution, puis les traités internationaux, la loi et les
  règlements, que le Code du travail réunit dans ses articles L et R, enfin,
  propres au droit du travail, les accords collectifs, interprofessionnels, de
  branche ou d'entreprise, et le contrat de travail. Une norme de rang
  inférieur doit en principe respecter les normes de rang supérieur ; c'est cet
  ordre que formalise l'ordre normatif des sources
  (\autoref{def:source}). En droit du travail, il ne suffit
  pourtant pas à régler un concours : une norme inférieure peut l'emporter sur
  une norme supérieure lorsqu'elle est plus favorable au salarié, ou lorsque la
  loi le lui permet.
 %%
\paragraph{Concours de normes.} Il y a concours lorsque plusieurs textes, tous
  applicables à la salariée à la date considérée, lui accordent un avantage
  ayant, selon la formule de la jurisprudence, « le même objet ou la même
  cause »~\cite{cassap1988} : ainsi l'indemnité de licenciement du Code du
  travail et celle de sa convention collective. En principe, ces avantages ne
  se cumulent pas : un seul est dû, et régler le concours, c'est dire lequel.
  En termes de programmation, une même fonction reçoit plusieurs définitions,
  de rangs différents, dont une seule doit être retenue pour un appel donné.
 %%
\paragraph{Ordre public et régime d'articulation.} Un concours ne se règle pas
  toujours de la même façon, car les sources du droit du travail ne
  s'articulent pas selon la seule hiérarchie des normes. La solution dépend du
  caractère d'\emph{ordre public} de la disposition de rang supérieur,
  c'est-à-dire de la latitude que celle-ci laisse aux sources de rang inférieur
  pour s'en écarter. Nous appelons \emph{régime d'articulation} la règle qui en
  résulte ; on en distingue trois~\cite{viepublique2021}. Sous l'ordre public
  \emph{absolu}, toute dérogation contractuelle ou conventionnelle est exclue,
  même plus favorable au salarié : la disposition supérieure s'applique seule.
  Sous l'ordre public \emph{social}, elle fixe un socle, que les sources
  inférieures peuvent améliorer : le concours se règle par le principe de
  faveur. Enfin, l'ordre public dérogatoire et les dispositions
  \emph{supplétives} autorisent des dérogations négociées dans un champ défini
  par le législateur : l'accord s'y applique, même moins favorable, et la
  disposition supérieure ne vaut qu'à défaut d'accord. Ces trois régimes
  évoquent respectivement une méthode que l'on ne peut redéfinir, une
  redéfinition dont le résultat n'est retenu que s'il est meilleur, et une
  implémentation par défaut.
  
Le régime dépend de la matière, et il vaut aussi entre accords collectifs de
niveaux différents : le calcul en cause doit donc être identifié avant que la
hiérarchie des sources puisse sélectionner une règle. Ainsi, l'article L2253-1
du \emph{Code du travail} dispose :
\begin{quote}\itshape
  Dans les matières énumérées au 1° à 13°, les stipulations de la convention de
  branche ou de l'accord couvrant un champ territorial ou professionnel plus
  large prévalent sur la convention d'entreprise conclue antérieurement ou
  postérieurement à la date de leur entrée en vigueur, sauf lorsque la
  convention d'entreprise assure des garanties au moins équivalentes. [\dots]
\end{quote}
Cette disposition illustre comment un texte peut empêcher qu'un calcul fourni
par une source de rang supérieur soit écarté par une source de rang inférieur
lorsque celle-ci produirait un résultat moins favorable. À l'inverse,
l'article L2253-3 dispose :
\begin{quote}\itshape
  Dans les matières autres que celles mentionnées aux articles L. 2253-1 et
  L. 2253-2, les stipulations de la convention d'entreprise conclue
  antérieurement ou postérieurement à la date d'entrée en vigueur de la
  convention de branche ou de l'accord couvrant un champ territorial ou
  professionnel plus large prévalent sur celles ayant le même objet prévues par
  la convention de branche ou l'accord couvrant un champ territorial ou
  professionnel plus large. En l'absence d'accord d'entreprise, la convention
  de branche ou l'accord couvrant un champ territorial ou professionnel plus
  large s'applique.
\end{quote}
Pour ces autres matières, l'accord d'entreprise, s'il existe, écarte l'accord
de branche ; à défaut, ce dernier demeure la règle de repli~\cite{frarticulation}.
Ce schéma de priorité est issu de la réforme de 2017, qui a renforcé la
négociation d'entreprise dans les champs fixés par la
loi~\cite{viepublique2017}.

\subsection{Question de recherche et contributions}\label{sec:question}
%%
Notre question de recherche est la suivante : \emph{comment un assistant
juridique peut-il organiser et résoudre des calculs concurrents de droit du
travail, identifier les informations nécessaires à leur sélection, et conserver
explicites et contrôlables les choix qui en résultent, afin d'aider un expert
juridique, éventuellement un grand modèle de langage (LLM), à construire une
plaidoirie sur une base maîtrisée et transparente ?} 

Nous décrivons, dans cet article, une modélisation des étapes du raisonnement juridique, esquissé \autoref{sec:juriste}, qui se trouve à la base de l'assistant \ALex~que nous développons (voir \autoref{fig:A-Lex}). Nous nous concentrons ici sur la définition de ces structures ainsi que sur la formalisation des étapes de raisonnement que ces structures supportent. L'implémentation concrète d'\ALex~n'est pas abordée, mais abstraite par une boucle d'interaction (\autoref{sec:ALex}) qui montre comment l'assistant utilise les structures introduites dans ce papier et que nous listons ci-dessous. Les ontologies juridiques décrites dans la \autoref{sec:vocabulaire} sont rappelées en caractères gras.
%%
\paragraph{Matières juridiques.}
Nous formalisons la notion de matière juridique comme un ensemble de types \emph{matières} équipés d'une relation de sous-typage (\autoref{def:matiere}). Une \emph{matière} $m$ étant un sous-type de $m'$ lorsque $m$ précise $m'$ (par exemple, \textsf{Prime} est un sous-type de \textsf{Montant}). Nous organisons les \emph{matières} dans un \emph{arbre de raffinement} (\autoref{def:raffinement}) qui permet à \ALex~de traduire une question juridique (une \emph{matière} abstraite) vers l'ensemble de ses concrétisations réalisables, étant donné un dossier client. Une \emph{matière} est concrète lorsqu'aucune autre ne la précise ; les matières concrètes sont les feuilles de l'\emph{arbre de raffinement}.
%%
\paragraph{Dossier client.}%%
Un dossier client est ici un \emph{dossier} formel (\autoref{def:dossier}), composé d'un ensemble de \emph{matières} et d'une valuation partielle qui associe à certaines d'entre elles une valeur ou un \emph{arbre normatif}. Une \emph{matière} du \emph{dossier} est \emph{établie} lorsqu'elle a une valeur, \emph{indéterminée} lorsqu'elle porte un arbre qu'il reste à résoudre, et \emph{ouverte} lorsqu'elle pose une question abstraite encore sans réponse. Un \emph{dossier} est créé lorsqu'une matière abstraite (la question posée) y est saisie. Les opérateurs de la \autoref{sec:op} permettent au juriste, en interaction avec \ALex, de conduire le \emph{dossier} vers un état où toutes ses \emph{matières} sont établies. 
%%
\paragraph{Sources juridiques.}
Comme pour les \emph{matières}, nous formalisons la notion de sources juridiques par un ensemble de types \emph{sources} équipés d'une relation de sous-typage (\autoref{def:source}). Une \emph{source} $s$ est un sous-type d'une source $s'$ lorsque $s$ précise $s'$ (par exemple, \textsf{ConventionBranche} est un sous-type d'\textsf{AccordCollectif}). Les \emph{sources} sont par ailleurs partiellement ordonnées selon la \textbf{hiérarchie des normes} en vigueur dans le droit du travail. Une \emph{source} est concrète lorsqu'elle est son unique sous-type.
%%
\paragraph{Arbres normatifs.} %%
Les \emph{arbres normatifs} (\autoref{def:normative-tree}) reposent sur un corpus~\cite{corpuslumo}, composé d'un ensemble de \emph{candidats} (\autoref{def:candidat}) typés par des \emph{matières} concrètes et indexés par leur \emph{source} concrète. L'implémentation Catala~\cite{merigoux2021} de ces fonctions n'est pas abordée dans ce papier, les \emph{candidats} ne retenant de ces fonctions que leur signature et une abstraction du calcul effectivement réalisé. Chaque \emph{candidat} est aussi équipé d'une \emph{politique d'arbitrage} (\autoref{def:arbitrage}) qui formalise son
\textbf{régime d'articulation}.  Un \emph{arbre normatif} structure des \emph{candidats} pour une même \emph{matière} $m$ en fonction de leur \emph{source} : la relation parent-enfant de l'arbre reflétant la \textbf{hiérarchie des normes}. Chaque arête de l'arbre est équipée d'un \emph{prédicat} (\autoref{def:applicability-predicate}) dont l'évaluation dépend du \emph{dossier} support et qui détermine si le candidat peut être retenu par le juriste, au titre d'une demande soutenue par ce \emph{dossier}. Deux candidats applicables d'un même \emph{arbre normatif} sont toujours en conflit et une norme doit l'emporter sur l'autre, deux candidats issus d'une même \emph{matière} abstraite typés par des \emph{matières} concrètes différentes peuvent parfois se cumuler. 

\paragraph{Étapes du raisonnement juridique assisté.}
Le raisonnement juridique est formalisé par quatre opérateurs sur les \emph{dossiers} (\autoref{sec:op}) : \emph{saisir} une \emph{matière}, fait ou question ; \emph{préciser} une \emph{matière} abstraite en la remplaçant, dans le même dossier, par celles de ses précisions que le dossier rend applicables ; \emph{élaguer} un \emph{arbre normatif} selon les faits établis ; \emph{résoudre} un arbre élagué en une valeur, issue du calcul du candidat que les \emph{politiques d'arbitrage} sélectionnent. Nous établissons que l'élagage possède une forme normale unique et que la résolution d'une matière abstraite vers un ensemble de valeurs possibles termine (\autoref{thm:boucle}). 

\begin{figure}[t]
\centering
\begin{tikzpicture}[
  scale=0.85, transform shape,
  x=1cm, y=1cm, font=\scriptsize,
  couche/.style={rotate=90, align=center, font=\scriptsize\bfseries},
  abs/.style={draw, rounded corners=2pt, align=center, inner sep=2.5pt,
    densely dashed, fill=white},
  conc/.style={draw, rounded corners=2pt, align=center, inner sep=2.5pt,
    thick, fill=white},
  cand/.style={draw, rounded corners=2pt, align=center, inner sep=2.5pt,
    fill=white},
  gris/.style={draw=elague, text=elague},
  vert/.style={draw=chemin, very thick},
  bleu/.style={draw=foret, very thick},
  racine/.style={circle, fill=foret, inner sep=1.8pt},
  module/.style={draw, align=center, inner sep=3pt, fill=white,
    font=\scriptsize\ttfamily},
  texte/.style={align=center, font=\tiny\itshape},
  arr/.style={-{Latex[length=1.6mm]}},
  arrv/.style={-{Latex[length=1.8mm]}, draw=chemin, very thick},
  arrb/.style={-{Latex[length=1.8mm]}, draw=foret, very thick},
  arrg/.style={-{Latex[length=1.6mm]}, draw=elague, densely dashed},
  lien/.style={-{Latex[length=1.6mm]}, densely dotted},
  lienb/.style={-{Latex[length=1.8mm]}, draw=foret, thick, densely dotted},
  pred/.style={font=\tiny, fill=bande, inner sep=1pt, align=center}
]
% Bandes
\fill[bande, rounded corners=3pt] (0,0.15) rectangle (13.6,-3.95);
\fill[bande, rounded corners=3pt] (0,-4.05) rectangle (13.6,-7.25);
\fill[bande, rounded corners=3pt] (0,-7.35) rectangle (13.6,-9.8);
\node[couche] at (0.35,-1.9) {2. Arbre de\\ raffinement};
\node[couche] at (0.35,-5.65) {1. Forêt\\ normative};
\node[couche] at (0.35,-8.57) {0. Corpus};

% Couche 2 : arbre de raffinement
\node[abs] (top) at (4.2,-0.35) {\(\Matter\)};
\node[abs] (dty) at (1.7,-1.1) {\(\Duration\)};
\node[abs] (mon) at (5.4,-1.1) {\(\Money\)};
\node[abs, vert, densely dashed] (ind) at (5.4,-1.85) {indemnité};
\node[abs, vert, densely dashed] (rup) at (6.9,-2.6) {indemnité de rupture};
\node[conc] (dur) at (1.7,-3.45) {\(\DureePreavis\)};
\node[conc, vert] (icp) at (3.9,-3.45) {\(\ICP\)};
\node[conc, vert] (lic) at (6.9,-3.45) {\(\Indem\)};
\node[draw=chemin, very thick, rounded corners=2pt, align=left, inner sep=3pt,
  fill=white] (q) at (11.2,-1.6)
  {\textit{question :} les indemnités\\ dues à la salariée};
\draw[arr] (top) -- (dty);
\draw[arr] (top) -- (mon);
\draw[arr] (dty) -- (dur);
\draw[arr] (mon) -- (ind);
\draw[arrv] (ind) -- node[pred, pos=0.5, left=1pt]{\(\PreavisNE\)} (icp);
\draw[arrv] (ind) -- (rup);
\draw[arrv] (q) -- (ind);
\draw[arrv] (rup) -- node[pred, right=2pt]{\(\ModeRupture=\Licenciement\)} (lic);

% Couche 1 : foret normative
\draw[fill=white] (1.7,-4.35) -- (1.2,-5.2) -- (2.2,-5.2) -- cycle;
\draw[draw=foret, very thick, fill=white] (3.9,-4.35) -- (3.4,-5.2) -- (4.4,-5.2) -- cycle;
\draw[lien] (dur) -- (1.7,-4.35);
\draw[arrb, densely dashed] (icp) -- (3.9,-4.35);
\node[racine] (r) at (6.9,-4.45) {};
\node[cand, gris] (old) at (5.6,-5.45) {Code, rédaction\\ antérieure};
\node[cand, bleu] (code) at (8.6,-5.45) {Code, R.~1234-2\\
  {\tiny\(\ConvBranche\mapsto\Social\)}};
\node[cand, gris] (s19) at (11.6,-4.5) {Syntec, art.~19};
\node[cand, gris] (hcr) at (11.6,-5.3) {HCR};
\node[cand, bleu] (s45) at (11.6,-6.6) {Syntec, art.~4.5};
\draw[arrb, densely dashed] (lic) -- (r);
\draw[arrg] (r) -- (old);
\draw[arrb] (r) -- node[pred, pos=0.45, right=3pt]{\(\DateT\geq\text{27/09/2017}\)} (code);
\draw[arrg] (code) -- (s19);
\draw[arrg] (code) -- (hcr);
\draw[arrb] (code) -- (s45);
\node[pred, anchor=west] at (10.55,-5.95)
  {\(\IDCC=1486\)\\ \(\DateT\geq\text{01/05/2023}\)};

% Couche 0 : corpus
\node[module, draw=foret, thick] (mpre) at (2.4,-8.15) {Loi\_preavis};
\node[module] (mhcr) at (5.6,-8.15) {Ccn\_1979\_\\ rupture\_du\_contrat};
\node[module, draw=foret, thick] (mlic) at (8.6,-8.15) {Loi\_indemnite\_\\ licenciement};
\node[module, draw=foret, thick] (msyn) at (11.6,-8.15) {Ccn\_1486\_\\ rupture\_du\_contrat};
\node[texte] (tpre) at (2.4,-9.25) {\S\ C. trav., L.~1234-1\\ et L.~1234-5};
\node[texte] (thcr) at (5.6,-9.25) {\S\ convention HCR\\ (IDCC 1979)};
\node[texte] (tlic) at (8.6,-9.25) {\S\ C. trav., L.~1234-9\\ et R.~1234-2};
\node[texte] (tsyn) at (11.6,-9.25) {\S\ convention Syntec,\\ art.~19 et 4.5};
\foreach \m/\t in {mpre/tpre, mhcr/thcr, mlic/tlic, msyn/tsyn}
  \draw[lien] (\m) -- (\t);
\draw[lienb] (code) -- (mlic);
\draw[lienb] (s45) -- (msyn);
\draw[lienb] (3.75,-5.2) -- (mpre);
\draw[lien] (1.7,-5.2) -- (mpre);
\end{tikzpicture}
\caption{Architecture d'\ALex{}. La couche~0, le corpus, réunit des modules
  Catala indépendants : chacun calcule la réponse d'un texte sans connaître les
  autres, et renvoie aux articles qu'il traduit (\S). La couche~1, la forêt
  normative, organise ces calculs en candidats, avec leurs conditions
  d'application et leurs politiques d'arbitrage, que le corpus laisse hors du
  code. La couche~2, l'arbre de raffinement, relie les questions abstraites aux
  matières concrètes, rangées sous leur type de valeurs
  (\autoref{fig:hierarchie}) ; les feuilles sont désignées par leur
  identifiant. En vert, le raffinement : le chemin part de la question de la
  salariée et se divise sous l'indemnité, entre \(\Indem\) et \(\ICP\), dues
  ensemble, qui ont chacune leur arbre normatif. En bleu, la forêt normative
  et son élagage : dans le premier arbre (racine en point), les candidats qui
  répondent au dossier fil rouge, avec les prédicats de leurs arêtes,
  jusqu'aux modules qui les calculent ; le second, esquissé, puise dans
  \texttt{Loi\_preavis} et dans le module Syntec. En gris, ce que le dossier
  écarte.}
\label{fig:A-Lex}
\end{figure}

\subsection{Exemple fil rouge}\label{sec:exemple}
Nous introduisons ici notre exemple fil rouge, tiré du corpus~\cite{corpuslumo}, qui illustrera, au fil des sections, les notions de notre formalisation à mesure qu'elles seront définies.

\textit{Une salariée d'une société de conseil relevant de la convention collective
Syntec (IDCC~1486) reçoit, en mars 2025, la notification de son licenciement,
qui n'est pas prononcé pour faute grave ; son employeur la dispense d'exécuter
son préavis. Ingénieure, elle relève de la catégorie des ingénieurs et cadres
de la convention, et non de celle des ETAM (employés, techniciens et agents de
maîtrise) ; elle compte dix ans d'ancienneté et son salaire de référence est
de 3\,000 euros. Elle demande les indemnités qui lui sont dues.}

Cette demande, « les indemnités qui lui sont dues », ne désigne encore aucun
calcul : c'est une \emph{matière abstraite}, au sens de la
\autoref{def:matiere}, que seuls les faits du dossier permettront de préciser.
Elle est, au départ, le seul contenu du dossier de la salariée
(\autoref{def:dossier}). Au fur et à mesure de l'exécution de cette question
initiale, on y ajoutera de nouvelles matières. La poursuite de l'exemple, au
fil de l'article, illustre le reste.

\subsection{Travaux connexes et positionnement}\label{sec:related}

Les travaux dits \emph{Rules as Code} étudient la représentation et
l'exécution disciplinées des dispositions juridiques~\cite{ali2025} :
LegalRuleML structure règles, sources, exceptions et
métadonnées~\cite{nazarenko2016,legalruleml}, L4 dote les contrats d'une
sémantique déontique et temporelle~\cite{watt2023}, Catala modélise les
règles par défaut et les exceptions de la loi~\cite{merigoux2021}. Nous
supposons ce problème résolu : les textes sont des objets formels qui en
contiennent la partie calculatoire, quel que soit le langage qui les
implémente, ici Catala.

Catala résout les définitions concurrentes d'une variable par une logique des
défauts à priorités~\cite{merigoux2021} : la sélection dépend de conditions
d'applicabilité et d'un ordre statique entre exceptions, non des valeurs
produites, alors que le principe de faveur dépend du résultat. Une exception
conditionnée à une comparaison de valeurs peut l'encoder, mais doit être
écrite dans le texte inférieur, pour chaque paire de sources, et la politique
se disperse dans les calculs ; notre corpus l'évite en isolant la comparaison
dans un module de concours. Les arbres normatifs poussent cette séparation
jusqu'au bout : calculs, sources et politiques d'arbitrage sont trois données
distinctes, et un même critère, la tolérance de chaque ancêtre le long du
chemin des sources, vaut pour tout régime. Catala structure les exceptions au
sein d'un texte, les arbres normatifs articulent des textes concurrents.

Les conflits entre normes sont classiquement traités par des règles
défaisables munies d'une relation de priorité, éventuellement dérivée par
des règles~\cite{prakkensartor1996,prakkensartor1997}, et par des logiques
défaisables où \emph{lex superior}, \emph{lex specialis} et \emph{lex
posterior} instancient la relation de supériorité~\cite{governatori2021},
les versions temporelles d'un texte étant des modifications de
normes~\cite{governatori2005}. Les arbres normatifs s'en distinguent sur
trois points : la priorité est une propriété de la norme supérieure, qui
énonce ce qu'elle tolère des sources inférieures, de sorte que
\(\Absolute\), \(\Suppletive\) et \(\Social\) coexistent sur un même
chemin ; la sélection sous \(\Social\) exige d'évaluer les deux candidats,
ce qu'une priorité statique n'exprime pas ; un fait manquant bloque la
résolution et engendre une demande, au lieu qu'une règle s'applique par
défaut. En contrepartie, la portée est plus étroite : candidats d'une seule
matière, frères mutuellement exclusifs plutôt qu'une relation générale de
spécificité, politiques non encore dérivées de textes tiers
(\autoref{sec:conclusion}). Enfin, les modèles du raisonnement juridique
comme construction et résolution
d'arguments~\cite{gordon1993pleadings,benchcapon2009argumentation} sont plus
larges : nous n'en automatisons que la part gouvernée par des calculs, des
conditions d'applicabilité et des politiques, les jugements interprétatifs
de l'expert étant représentés par les saisies (\autoref{sec:saisir}).
